\documentclass[12pt,a4paper]{article}
\usepackage[utf8]{inputenc}
\usepackage[vietnamese]{babel}
\usepackage{amsmath}
\usepackage{amssymb}
\usepackage{graphicx}
\usepackage{xcolor}
\usepackage{listings}
\usepackage{hyperref}
\usepackage{geometry}
\usepackage{fancyhdr}
\usepackage{multirow}
\usepackage{array}
\usepackage{tabularx}
\usepackage{float}

\geometry{top=2.5cm, bottom=2.5cm, left=2.5cm, right=2.5cm}
\setlength{\parindent}{1.5cm}
\setlength{\parskip}{0.3cm}

% Code style
\lstset{
    basicstyle=\ttfamily\small,
    breaklines=true,
    keywordstyle=\color{blue},
    commentstyle=\color{gray},
    stringstyle=\color{red},
    showstringspaces=false,
    tabsize=2,
    language=Python,
    frame=single,
    numbers=left,
    numberstyle=\tiny,
    backgroundcolor=\color{gray!10}
}

% Header/Footer
\pagestyle{fancy}
\fancyhf{}
\lhead{\small \textit{Báo cáo Tiểu luận}}
\rhead{\small \thepage}

% No title page - start directly with TOC as requested
\begin{document}

{\Large \textbf{NỀN TẢNG Y TẾ ĐIỆN TỬ TÍCH HỢP TRỌNG TÂM AI}}

{\Large \textit{Health-Micro-AI: Kiến Trúc Microservices và Ứng Dụng Trí Tuệ Nhân Tạo}}

\vspace{2cm}
\tableofcontents
\newpage

% ══════════════════════════════════════════════════════════════════════════════
% CHƯƠNG 1: KIẾN TRÚC MICROSERVICES
% ══════════════════════════════════════════════════════════════════════════════

\section{Kiến Trúc Microservices và Domain-Driven Design}

\subsection{Từ Monolithic đến Microservices}

Kiến trúc phần mềm truyền thống (Monolithic Architecture) gom tất cả chức năng vào một ứng dụng duy nhất, khó bảo trì, khó mở rộng. Khi hệ thống có hàng trăm nghiệp vụ và hàng tỉ giao tác, một số vấn đề nổi cộm:

\begin{itemize}
    \item \textbf{Kết dính chặt chẽ (Tight Coupling)}: Thay đổi một tính năng dễ ảnh hưởng toàn hệ.
    \item \textbf{Khó mở rộng (Poor Scalability)}: Phải duplicate toàn bộ ứng dụng.
    \item \textbf{Khó triển khai (Deployment Challenge)}: Một lỗi nhỏ làm crash toàn bộ.
    \item \textbf{Khó bảo trì (Maintenance Burden)}: Codebase khổng lồ, khó hiểu.
\end{itemize}

\textbf{Giải pháp:} Kiến trúc Microservices (vi dịch vụ) tách ứng dụng thành các dịch vụ nhỏ, độc lập, giao tiếp qua API. Mỗi dịch vụ:
\begin{itemize}
    \item Quản lý một phần nghiệp vụ riêng biệt (Single Responsibility Principle).
    \item Có cơ sở dữ liệu riêng (Database Per Service Pattern).
    \item Phi tập trung, có thể triển khai độc lập.
    \item Cho phép sử dụng công nghệ khác nhau (polyglot architecture).
\end{itemize}

\textbf{Ưu điểm của Microservices:}
\begin{enumerate}
    \item \textbf{Khả năng mở rộng cao} – Scale từng service riêng lẻ theo load.
    \item \textbf{Khả năng phục hồi} – Một service lỗi không làm crash toàn hệ.
    \item \textbf{Linh hoạt công nghệ} – Mỗi service chọn tech stack tối ưu.
    \item \textbf{Dễ bảo trì} – Team nhỏ quản lý một service cụ thể.
    \item \textbf{Tốc độ phát triển} – Deploy nhanh mà không cần release toàn hệ.
\end{enumerate}

\textbf{Thách thức của Microservices:}
\begin{enumerate}
    \item \textbf{Độ phức tạp giao tiếp} – Phải quản lý distributed calls, timeout, retry.
    \item \textbf{Tính nhất quán dữ liệu} – Phải kết hợp dữ liệu từ nhiều database.
    \item \textbf{Logging và Monitoring} – Cần centralized tracing để debug.
    \item \textbf{Deployment} – Cần orchestration tool (Docker, Kubernetes).
\end{enumerate}

\subsection{Thiết Kế Microservices Theo Domain-Driven Design}

\textbf{Domain-Driven Design (DDD)} là phương pháp thiết kế phần mềm bằng cách mô hình hóa các nghiệp vụ thực tế (domain). Thay vì tập trung vào kỹ thuật (framework, database), DDD tập trung vào \textit{ngôn ngữ chung} giữa kỹ sư và chuyên gia kinh doanh.

\textbf{Các khái niệm cốt lõi của DDD:}

\begin{description}
    \item[Ubiquitous Language (Ngôn ngữ chung)] Dùng từ vựng thống nhất để nói về domain trong code, documentation, và hội thoại.
    
    \item[Bounded Context (Ngữ cảnh giới hạn)] Chia domain lớn thành các phần nhỏ, mỗi phần có:
    \begin{itemize}
        \item Ngôn ngữ riêng
        \item Model riêng
        \item Ranh giới rõ ràng
        \item Giao tiếp qua API/Events
    \end{itemize}
    
    \item[Aggregate (Tập hợp)] Nhóm các entities liên quan thành một đơn vị quản lý trans actional. Chỉ có thể modify thông qua aggregate root.
    
    \item[Entity (Thực thể)] Đối tượng có identity không đổi. Ví dụ: Patient (bệnh nhân) có ID không đổi ngay cả khi thay đổi tên/địa chỉ.
    
    \item[Value Object (Đối tượng giá trị)] Không có identity riêng, chỉ được định nghĩa bởi giá trị. Ví dụ: Amount (số tiền), Priority (mức độ ưu tiên). Không thể thay đổi – phải tạo mới.
    
    \item[Domain Event (Sự kiện miền)] Ghi lại sự kiện quan trọng xảy ra trong domain. Các dịch vụ khác lắng nghe để sync dữ liệu. Ví dụ: \textit{PatientRegistered}, \textit{PrescriptionSubmitted}.
    
    \item[Repository (Kho lưu trữ)] Abstraction để truy cập data. Che dấu chi tiết về cơ sở dữ liệu (SQL, NoSQL, v.v.).
    
    \item[Domain Service (Dịch vụ miền)] Logic nghiệp vụ không thuộc entity/value object. Ví dụ: RecommendationEngine (tính toán gợi ý).
\end{description}

\textbf{Kiến trúc 4-Layer DDD:}

Mỗi Bounded Context (microservice) được tổ chức thành 4 lớp:

\begin{enumerate}
    \item \textbf{Interface Layer (Lớp giao diện)} – REST API, HTTP endpoints, nhận request từ client.
    
    \item \textbf{Application Layer (Lớp ứng dụng)} – CQRS (Commands \& Queries). Command (viết), Query (đọc). Orchestrate domain logic.
    
    \item \textbf{Domain Layer (Lớp miền)} – Aggregates, Entities, Value Objects, Domain Services, Domain Events. Chứa logic nghiệp vụ.
    
    \item \textbf{Infrastructure Layer (Lớp cơ sở hạ tầng)} – Database, Repository Implementation, External Service Calls.
\end{enumerate}

\begin{verbatim}
******* DDD 4-LAYER ARCHITECTURE *******

[1] INTERFACE LAYER (REST API)
    └─> GET /patients, POST /consultations
    
[2] APPLICATION LAYER (CQRS Pattern)
    ├─ Commands: CreatePatientCommand
    └─ Queries: ListPatientsQuery
    
[3] DOMAIN LAYER (Business Logic)
    ├─ Patient Aggregate (root entity)
    └─ Domain Events: PatientCreated
    
[4] INFRASTRUCTURE LAYER (Database)
    ├─ PatientRepository
    └─ PostgreSQL + SQLite ORM
    
*****  Data flows from top to bottom  *****
\end{verbatim}

Lợi ích của kiến trúc này:
\begin{itemize}
    \item \textbf{Độc lập kiểm thử} – Test domain logic không cần database.
    \item \textbf{Dễ bảo trì} – Mỗi lớp có trách nhiệm rõ.
    \item \textbf{Linh hoạt công nghệ} – Swap database dễ dàng.
    \item \textbf{Reusable} – Domain logic không phụ thuộc framework.
\end{itemize}

\subsection{Case Study: Hệ Thống E-Commerce Healthcare}

\textbf{Bài toán đặt vấn đề:}

Trong bối cảnh phát triển của ngành y tế điện tử (E-health), người dùng Việt Nam gặp những khó khăn:

\begin{enumerate}
    \item \textbf{Thông tin sức khỏe phân tán} – Bệnh nhân phải gửi lại hồ sơ ở mỗi cơ sở y tế.
    
    \item \textbf{Tìm kiếm tư vấn y tế khó} – Thông tin trên Internet rối ren, không chuyên môn, dễ gây hoang mang.
    
    \item \textbf{Tra cứu thuốc phức tạp} – Không biết thuốc nào phù hợp, tương tác với thuốc khác.
    
    \item \textbf{Gợi ý cạnh tranh} – Khó biết sản phẩm nào phù hợp với tình trạng sức khỏe riêng.
    
    \item \textbf{Đơn thuốc không được quản lý} – Không có lịch sử đơn thuốc, không cảnh báo tương tác.
\end{enumerate}

\textbf{Giải pháp: Nền tảng Health-Micro-AI}

Dự án này xây dựng một nền tảng e-commerce y tế hoàn chỉnh với:

\begin{itemize}
    \item \textbf{Quản lý bệnh nhân tập trung} – Hồ sơ y tế điện tử (EHR – Electronic Health Record).
    \item \textbf{Chatbot AI tư vấn y tế} – Sử dụng GraphRAG (Neo4j + FAISS + LLM Gemini) để trả lời câu hỏi y tế dựa trên kiến thức khoa học.
    \item \textbf{Gợi ý cá nhân hóa} – Sử dụng Collaborative Filtering và GNN để gợi ý sản phẩm phù hợp.
    \item \textbf{Quản lý đơn thuốc} – Tạo, lưu giữ, và theo dõi đơn thuốc.
    \item \textbf{Danh mục y khoa} – Theo ICD-10, phân loại bệnh tật và sản phẩm.
    \item \textbf{Đánh giá cộng đồng} – Người dùng đánh giá hiệu quả điều trị, chia sẻ kinh nghiệm.
\end{itemize}

Bằng cách sử dụng Microservices + DDD, hệ thống này:
\begin{itemize}
    \item Cho phép mở rộng riêng từng chức năng (ví dụ: scale AI service khi tầng lạnh tăng).
    \item Bảo vệ từng miền nghiệp vụ bằng Bounded Context.
    \item Giảm độ phức tạp và tăng khả năng bảo trì.
\end{itemize}

% ══════════════════════════════════════════════════════════════════════════════
% CHƯƠNG 2: HỆ THỐNG E-COMMERCE HEALTHCARE
% ══════════════════════════════════════════════════════════════════════════════

\newpage
\section{Hệ Thống E-Commerce Healthcare: Yêu Cầu và Phân Rã}

\subsection{Xác Định Yêu Cầu Chức Năng}

\textbf{Yêu cầu chức năng (Functional Requirements):}

Hệ thống phải hỗ trợ các vai trò:
\begin{enumerate}
    \item \textbf{Bệnh nhân (Customer)} – Tìm kiếm, mua thuốc, tư vấn AI, quản lý đơn thuốc.
    
    \item \textbf{Quản trị viên (Admin)} – Quản lý sản phẩm, danh mục, người dùng, duyệt đơn thuốc.
    
    \item \textbf{Nhân viên (Staff)} – Xử lý đơn thuốc, kiểm kho, cập nhật tồn kho.
\end{enumerate}

\textbf{Chức năng nòng cốt:}
\begin{itemize}
    \item \textbf{Xác thực và phân quyền} – Đăng ký, đăng nhập (JWT), phân vai trò.
    
    \item \textbf{Quản lý bệnh nhân} – Hồ sơ y tế, thông tin cá nhân, lịch khám.
    
    \item \textbf{Quản lý sản phẩm} – Danh mục thuốc, tìm kiếm, filter theo triệu chứng.
    
    \item \textbf{Chatbot AI tư vấn} – Trả lời câu hỏi y tế, gợi ý sản phẩm.
    
    \item \textbf{Gợi ý cá nhân hóa} – Recommend sản phẩm dựa trên hành vi và sở thích.
    
    \item \textbf{Quản lý đơn thuốc} – Tạo, lưu trữ, theo dõi trạng thái đơn.
    
    \item \textbf{Đánh giá và bình luận} – Người dùng đánh giá sản phẩm, chia sẻ kinh nghiệm.
\end{itemize}

\textbf{Yêu cầu phi chức năng (Non-Functional Requirements):}
\begin{itemize}
    \item \textbf{Khả năng mở rộng} – Hỗ trợ hàng triệu người dùng, mười triệu giao tác/ngày.
    
    \item \textbf{Độ tin cậy} – Uptime $\geq 99.9\%$ (SLA 3-nines).
    
    \item \textbf{Hiệu suất} – Response time $< 200$ms cho hầu hết request.
    
    \item \textbf{Bảo mật} – Bảo vệ thông tin bệnh nhân (HIPA, GDPR-like compliance).
    
    \item \textbf{Tính nhất quán dữ liệu} – Sử dụng event-driven architecture để đảm bảo consistency.
\end{itemize}

\subsection{Phân Rã Hệ Thống Theo DDD}

Bằng cách phân tích domain, ta xác định được các \textbf{Bounded Contexts} (ngữ cảnh giới hạn):

\begin{enumerate}
    \item \textbf{User Management Context} – Quản lý tài khoản, xác thực, phân quyền.
    
    \item \textbf{Patient Management Context} – Hồ sơ y tế, thông tin cá nhân bệnh nhân.
    
    \item \textbf{Product Catalog Context} – Danh mục sản phẩm, tìm kiếm, phân loại.
    
    \item \textbf{Prescription Management Context} – Đơn thuốc, giỏ hàng tạm thời.
    
    \item \textbf{AI Consulting Context} – Chatbot GraphRAG, tư vấn y tế.
    
    \item \textbf{Recommendation Context} – Gợi ý sản phẩm, machine learning.
    
    \item \textbf{Review \& Rating Context} – Đánh giá, bình luận sản phẩm.
\end{enumerate}

\textbf{Context Map (Bản đồ giao tiếp giữa Contexts):}

\begin{verbatim}
========== BOUNDED CONTEXTS & SERVICE TOPOLOGY ==========

ENTRY POINT
  [API GATEWAY :8000]
     Authentication, Routing, Rate Limiting
     

IDENTITY & ACCESS
  [AUTH SERVICE :8012]
     * User Registration
     * JWT Token Management  
     * Role-Based Access Control


USER & HEALTH RECORDS
  [PATIENT SERVICE :8001]
     * Patient Profiles
     * Medical History
     * Insurance Info
     

PRODUCT CATALOG
  [PHARMACY SERVICE :8002]
     * Drug Repository
     * Batch Tracking
     * Inventory Management
     

MEDICAL REFERENCE
  [MEDICAL-CATALOG SERVICE :8003]
     * ICD-10 Classification
     * Disease-Product Mapping
     * Treatment Protocols


PRESCRIPTION MANAGEMENT
  [PRESCRIPTION SERVICE :8004]
     * Prescription Workflow
     * Drug Interaction Check
     * Cart Management


LOGISTICS & FULFILLMENT
  [DISPENSING SERVICE :8005]
     * Order Fulfillment
     * Payment Processing
     * Shipping Coordination


COMMUNITY FEEDBACK
  [MEDICAL-REVIEW SERVICE :8006]
     * Rating System
     * User Reviews
     * Reputation Store


AI ADVISORY
  [CLINICAL-ADVISORY SERVICE :8013]  
     * GraphRAG Engine
     * Health Consultation
     * Product Recommendations


PERSONALIZATION ENGINE
  [RECOMMENDER-AI SERVICE :8009]
     * Collaborative Filtering
     * Behavior Analysis
     * Preference Learning


EVENT BUS (RabbitMQ)
  Domain Events: PrescriptionSubmitted, InventoryUpdated
  Publish-Subscribe Model for Eventual Consistency

========================================================
\end{verbatim}

\subsection{Công Nghệ và Stack}

\textbf{Backend Services:}
\begin{itemize}
    \item \textbf{Framework:} Django + Django REST Framework (chính), FastAPI (AI services).
    \item \textbf{Database:} PostgreSQL (tất cả services), SQLite (API Gateway).
    \item \textbf{AI/ML:} Neo4j (Knowledge Graph), FAISS (Vector Store), Gemini LLM (Cloud-based).
    \item \textbf{Message Broker:} RabbitMQ (event distribution).
    \item \textbf{Caching:} Redis (optional, cho production).
\end{itemize}

\textbf{Deployment:}
\begin{itemize}
    \item \textbf{Containerization:} Docker, Docker Compose.
    \item \textbf{Orchestration:} Manual setup (Docker Compose) hoặc Kubernetes (future).
\end{itemize}

\textbf{Development:}
\begin{itemize}
    \item \textbf{Language:} Python 3.9+.
    \item \textbf{Version Control:} Git.
    \item \textbf{Testing:} pytest, unittest.
\end{itemize}

\subsection{Dữ Liệu Khởi Tạo Thực Tế (Seed Data)}

Để chứng minh tính thực tiễn và khả năng đáp ứng của hệ thống, cơ sở dữ liệu đã được khởi tạo sẵn \textbf{10 mặt hàng y tế} mang tính đại diện cao (bao gồm cả thuốc kê đơn và thuốc không kê đơn - OTC), tập trung vào nhóm bệnh tiêu hóa và kháng sinh. Danh sách 10 mặt hàng đang được kinh doanh và quản lý trên hệ thống E-commerce bao gồm:

\begin{table}[H]
    \centering
    \begin{tabularx}{\textwidth}{|c|X|l|l|l|}
        \hline
        \textbf{STT} & \textbf{Tên Sản Phẩm} & \textbf{Hoạt Chất} & \textbf{Danh Mục} & \textbf{Trạng Thái} \\
        \hline
        1 & Phosphalugel 20\% gel & Aluminum Phosphate & Tiêu hóa (Antacid) & OTC \\
        \hline
        2 & Nexium 20mg & Esomeprazole & Tiêu hóa (PPI) & Kê đơn \\
        \hline
        3 & Omeprazole 20mg & Omeprazole & Tiêu hóa (PPI) & Kê đơn \\
        \hline
        4 & Ranitidine 150mg & Ranitidine HCl & Tiêu hóa (H2 Blocker) & Kê đơn \\
        \hline
        5 & Amoxicillin 500mg & Amoxicillin & Kháng sinh & Kê đơn \\
        \hline
        6 & Metronidazole 500mg & Metronidazole & Kháng sinh & Kê đơn \\
        \hline
        7 & Smecta 3g & Diosmectite & Tiêu hóa (Mucosal) & OTC \\
        \hline
        8 & Domperidone 10mg & Domperidone & Tiêu hóa (Antiemetic) & OTC \\
        \hline
        9 & Lacidofil (Probiotic) & Lactobacillus rhamnosus & Men vi sinh & OTC \\
        \hline
        10 & Duspatalin 200mg & Mebeverine HCl & Tiêu hóa (Chống co thắt) & OTC \\
        \hline
    \end{tabularx}
    \caption{Danh sách 10 mặt hàng y tế được quản lý trong hệ thống}
    \label{tab:10products}
\end{table}

Các mặt hàng này được liên kết trực tiếp với AI Context (GraphRAG) để chatbot có thể đưa ra tư vấn và trỏ thẳng người dùng đến \textit{Trang Chi tiết Sản phẩm} nhằm bỏ vào giỏ hàng.

% ══════════════════════════════════════════════════════════════════════════════
% CHƯƠNG 3: THIẾT KỾ CHI TIẾT
% ══════════════════════════════════════════════════════════════════════════════

\newpage
\section{Thiết Kế Chi Tiết Các Services}

\subsection{Kiến Trúc Tổng Thể Hệ Thống}

\begin{figure}[H]
    \centering
    \begin{verbatim}
============= HEALTH-MICRO-AI ARCHITECTURE =============

CLIENT LAYER
  [Browser] http://localhost:8000

GATEWAY LAYER (Port 8000)
  [API GATEWAY] Django + SQLite
    * Render Templates
    * Route & Auth
    * Session Mgmt

SERVICE LAYER (Microservices)
  PORT    SERVICE              FRAMEWORK    DB
  ----    -------              ---------    --
  8012    Auth Service         FastAPI      PostgreSQL
  8001    Patient Service      Django       PostgreSQL
  8002    Pharmacy Service     Django       PostgreSQL
  8004    Prescription Svc     FastAPI      PostgreSQL
  8003    Medical-Catalog      Django       PostgreSQL
  8005    Dispensing Service   Django       PostgreSQL  
  8006    Medical-Review       FastAPI      PostgreSQL
  8013    Clinical-Advisory    FastAPI      Neo4j
  8009    Recommender-AI       Django       Neo4j

INFRASTRUCTURE LAYER
  [PostgreSQL 5432]    [Neo4j 7687]    [RabbitMQ]
  * patient_db         * Knowledge     * Event Bus
  * pharmacy_db          Graph          * Topics
  * prescription_db    * User-Item
  * medical_catalog      Relationships

EVENT-DRIVEN COMMUNICATION (Async)
  Services publish Domain Events -> RabbitMQ -> Subscribers

============================================================
    \end{verbatim}
    \caption{Kiến trúc toàn cục hệ thống Health-Micro-AI (11 services + infrastructure)}
    \label{fig:architecture}
\end{figure}

\subsection{Từng Service – Một Trong Các Bounded Contexts}

\subsubsection{1. Auth Service (Port :8012, FastAPI)}

\textbf{Mục tiêu:} Quản lý xác thực (authentication) và phân quyền (authorization).

\textbf{Chức năng:}
\begin{itemize}
    \item Đăng ký tài khoản mới.
    \item Đăng nhập, cấp JWT token (access\_token, refresh\_token).
    \item Xác thực token (verify, refresh).
    \item Quản lý vai trò (Role), quyền (Permission).
    \item Quản lý user (lock account, change password).
\end{itemize}

\textbf{Domain Model (DDD):}
\begin{itemize}
    \item \textbf{Entities:} User (aggregate root), Role, Permission.
    \item \textbf{Value Objects:} UserId, Email, CredentialHash, Token, JWTClaims.
    \item \textbf{Domain Events:} UserRegistered, UserLoggedIn, TokenGenerated, PasswordChanged, RoleAssigned.
    \item \textbf{Domain Services:} AuthenticationService, AuthorizationService, TokenManagementService.
\end{itemize}

\textbf{Endpoints chính:}
\begin{verbatim}
POST   /api/auth/register          – Đăng ký
POST   /api/auth/login             – Đăng nhập
POST   /api/auth/validate          – Xác thực token
POST   /api/auth/refresh           – Refresh token
GET    /api/roles                  – Lấy danh sách vai trò
POST   /api/permissions/{role}     – Gán quyền cho vai trò
\end{verbatim}

\textbf{Database PostgreSQL:}
\begin{verbatim}
auth_db
+-- users
|   +-- id, email, password_hash, fullname, created_at, is_active
+-- roles
|   +-- id, name, description
+-- permissions
|   +-- id, name, description
+-- user_roles
    +-- user_id, role_id
\end{verbatim}

\subsubsection{2. Patient Service (Port :8001, Django)}

\textbf{Mục tiêu:} Quản lý thông tin bệnh nhân, hồ sơ y tế.

\textbf{Chức năng:}
\begin{itemize}
    \item Tạo and cập nhật hồ sơ bệnh nhân (tên, tuổi, giới tính, địa chỉ).
    \item Lưu trữ thông tin bảo hiểm y tế (Insurance Code).
    \item Theo dõi lịch sử y tế (medical history, allergies, chronic diseases).
    \item Compute membership tier dựa trên total spending.
\end{itemize}

\textbf{Domain Model:}
\begin{itemize}
    \item \textbf{Entities:} Patient (aggregate root), MedicalHistory, Allergy.
    \item \textbf{Value Objects:} PatientId, FullName, DateOfBirth, InsuranceCode, BloodType.
    \item \textbf{Domain Events:} PatientRegistered, MedicalHistoryUpdated, AllergiesAdded.
\end{itemize}

\textbf{Database PostgreSQL:}
\begin{verbatim}
patient_db
+-- patients
|   +-- id, user_id, fullname, dob, gender, phone, email,
|       address, insurance_code, blood_type, total_spent,
|       membership_tier, created_at, updated_at
+-- medical_histories
|   +-- id, patient_id, condition, diagnosed_date, status
+-- allergies
    +-- id, patient_id, allergen, reaction_severity
\end{verbatim}

\subsubsection{3. Pharmacy Service (Port :8002, Django)}

\textbf{Mục tiêu:} Quản lý danh mục sản phẩm, tồn kho, batch tracking.

\textbf{Chức năng:}
\begin{itemize}
    \item Quản lý sản phẩm (thuốc, thiết bị y tế) – thêm, xóa, cập nhật.
    \item Theo dõi tồn kho (inventory) theo batch/lô, hạn sử dụng (expiry date).
    \item Tìm kiếm sản phẩm theo triệu chứng (symptom\_tags).
    \item Filter theo danh mục, giá, nhà sản xuất.
\end{itemize}

\textbf{Domain Model:}
\begin{itemize}
    \item \textbf{Entities:} Product (aggregate root), Batch.
    \item \textbf{Value Objects:} ProductId, ProductCode, DosageStrength, Price, StockQuantity.
    \item \textbf{Domain Events:} ProductCreated, BatchReceived, InventoryAdjusted.
    \item \textbf{Domain Services:} InventoryService, ProductSearchService.
\end{itemize}

\textbf{Database PostgreSQL:}
\begin{verbatim}
pharmacy_db
+-- products
|   +-- id, product_code, name, generic_name, category,
|       dosage_form, strength, price, stock_qty,
|       requires_prescription, manufacturer, origin, 
|       symptom_tags (JSON array), description
+-- batches
|   +-- id, product_id, batch_number, qty, mfg_date,
|       expiry_date, location
+-- inventory_locations
    +-- id, product_id, location_name, qty, reserved_qty
\end{verbatim}

\subsubsection{4. Prescription Service (Port :8004, FastAPI)}

\textbf{Mục tiêu:} Quản lý đơn thuốc, giỏ hàng (cart) – tạm thời, fulfillment workflow.

\textbf{Chức năng:}
\begin{itemize}
    \item Tạo đơn thuốc (từ bác sĩ hoặc tự tạo).
    \item Thêm/xóa sản phẩm vào đơn (cart management).
    \item Submit đơn để xử lý (trigger Dispensing Service).
    \item Lưu trữ lịch sử đơn thuốc.
    \item Kiểm tra tương tác thuốc (drug interaction check).
\end{itemize}

\textbf{Domain Model:}
\begin{itemize}
    \item \textbf{Entities:} Prescription (aggregate root), PrescriptionLineItem (child entity).
    \item \textbf{Value Objects:} PrescriptionId, CustomerId, ProductId, Quantity, Dosage, DrugRequirement.
    \item \textbf{Domain Events:} PrescriptionCreated, ItemAdded, ItemRemoved, Submitted, Confirmed, Fulfilled, Cancelled.
    \item \textbf{Domain Services:} PrescriptionCreationService, CartManagementService, DrugInteractionService.
\end{itemize}

\textbf{Endpoints chính:}
\begin{verbatim}
POST   /api/prescriptions              – Tạo đơn thuốc
POST   /api/prescriptions/{id}/items   – Thêm sản phẩm
DELETE /api/prescriptions/{id}/items/{item_id} – Xóa
PUT    /api/prescriptions/{id}/submit  – Submit đơn
GET    /api/prescriptions/{id}         – Lấy chi tiết
GET    /api/prescriptions/customer/{cust_id} – Lịch sử
\end{verbatim}

\subsubsection{5. Medical-Catalog Service (Port :8003, Django)}

\textbf{Mục tiêu:} Danh mục y khoa – phân loại bệnh tật (ICD-10), sản phẩm.

\textbf{Chức năng:}
\begin{itemize}
    \item Danh sách bệnh tật theo ICD-10.
    \item Phân loại sản phẩm theo mục đích điều trị.
    \item Mô tả triệu chứng, chẩn đoán.
\end{itemize}

\textbf{Endpoints:}
\begin{verbatim}
GET /api/diseases?search=gastric     – Tìm bệnh
GET /api/diseases/{icd_code}         – Chi tiết bệnh
GET /api/categories                  – Danh sách danh mục
GET /api/categories/{cat_id}/products – Sản phẩm danh mục
\end{verbatim}

\subsubsection{6. Medical-Review Service (Port :8006, FastAPI)}

\textbf{Mục tiêu:} Đánh giá và bình luận sản phẩm từ người dùng.

\textbf{Chức năng:}
\begin{itemize}
    \item Thêm bình luận/đánh giá sản phẩm (5-star rating).
    \item Xác thực review từ confirmed purchases.
    \item Auto-aggregate rating trung bình.
    \item Upvote helpful reviews.
\end{itemize}

\textbf{Endpoints:}
\begin{verbatim}
POST /api/reviews                    – Thêm review
GET  /api/products/{prod_id}/reviews – Lấy reviews sản phẩm
POST /api/reviews/{id}/upvote        – Upvote review
\end{verbatim}

\subsubsection{7. Clinical-Advisory Service (Port :8013, FastAPI + GraphRAG)}

\textbf{Mục tiêu:} Chatbot AI tư vấn y tế (sẽ chi tiết ở Chương 4).

\subsubsection{8. Recommender-AI Service (Port :8009, Django)}

\textbf{Mục tiêu:} Gợi ý cá nhân hóa (sẽ chi tiết ở Chương 4).

\subsubsection{9. API Gateway (Port :8000, Django)}

\textbf{Mục tiêu:} Unified entry point – routing, authentication, rate limiting.

\textbf{Chức năng:}
\begin{itemize}
    \item Nơi duy nhất để frontend gọi (tránh CORS issues).
    \item Render HTML templates (web UI).
    \item Route request đến microservice phù hợp.
    \item Xác thực JWT token từ header.
    \item Rate limiting (mencegah abuse).
    \item Logging, correlation ID.
\end{itemize}

\textbf{Endpoints chính:}
\begin{verbatim}
GET  /                               – Trang chủ
GET  /health                         – Health check
POST /api/auth/login                 – Đăng nhập (proxy to Auth)
GET  /api/patients/{id}              – Proxy to Patient Service
POST /api/products/search            – Proxy to Pharmacy
GET  /api/chat                       – Web UI chat
\end{verbatim}

\subsection{Kết Nối Giữa Các Services}

Services giao tiếp qua 2 cách:

\begin{enumerate}
    \item \textbf{Synchronous (HTTP/REST):} Yêu cầu & phản hồi trực tiếp.
    \begin{itemize}
        \item Ví dụ: Prescription Service gọi Pharmacy Service để kiểm tra tồn kho.
        \item Ưu điểm: Đơn giản, immediate feedback.
        \item Nhược điểm: Coupling chặt chẽ, lỗi một service ảnh hưởng service khác.
    \end{itemize}
    
    \item \textbf{Asynchronous (Message Broker – RabbitMQ):} Publish-Subscribe via domain events.
    \begin{itemize}
        \item Ví dụ: Dispensing Service publish ``OrderFulfilled'' event, Pharmacy Service lắng nghe để update tồn kho.
        \item Ưu điểm: Loose coupling, resilient, scalable.
        \item Nhược điểm: Eventual consistency (dữ liệu không ngay lập tức nhất quán).
    \end{itemize}
\end{enumerate}

\textbf{Event Flow Example:}

\begin{verbatim}
1. Customer tạo Prescription
   → Prescription Service tạo PrescriptionCreated event

2. Event đẩy vào RabbitMQ Topic Exchange (prescription.events)
   
3. Các service subscribe:
   • Pharmacy Service: Kiểm tra tồn kho
   • Clinical-Advisory Service: Update user behavior graph
   • Recommendation Service: Ghi nhận interaction

4. Mỗi service xử lý event theo cách riêng (async, không block)
   
5. Nếu có lỗi, RabbitMQ retry theo exponential backoff
\end{verbatim}

% ══════════════════════════════════════════════════════════════════════════════
% CHƯƠNG 4: DỊCH VỤ AI
% ══════════════════════════════════════════════════════════════════════════════

\newpage
\section{Dịch Vụ AI: Chatbot GraphRAG và Gợi Ý Cá Nhân Hóa}

\subsection{Cơ Sở Lý Thuyết}

\subsubsection{Retrieval-Augmented Generation (RAG)}

\textbf{Định nghĩa:} RAG là kĩ thuật kết hợp \textit{Information Retrieval} (tìm kiếm thông tin) với \textit{Generative AI} (sinh văn bản). Thay vì chỉ dùng LLM để sinh câu trả lời (có nguy cơ hallucination – tạo ra thông tin sai), RAG:

\begin{enumerate}
    \item Truy xuất (retrieve) các tài liệu liên quan từ knowledge base.
    \item Tăng cường (augment) prompt bằng các tài liệu này.
    \item Yêu cầu LLM sinh câu trả lời dựa trên tài liệu (grounded).
\end{enumerate}

\textbf{Công thức toán:}
$$P(\text{answer}|\text{question}, \text{retrieved\_docs}) \propto P(\text{docs}|\text{question}) \cdot P(\text{answer}|\text{docs}, \text{question})$$

Trong đó:
\begin{itemize}
    \item $P(\text{docs}|\text{question})$ là xác suất tài liệu liên quan (retrieval stage).
    \item $P(\text{answer}|\text{docs}, \text{question})$ là xác suất câu trả lời cho trước tài liệu (generation stage).
\end{itemize}

\textbf{Lợi ích của RAG:}
\begin{itemize}
    \item \textbf{Giảm hallucination} – LLM phải trích dẫn từ tài liệu thực tế.
    \item \textbf{Dễ cập nhật} – Chỉ cần thêm tài liệu mới, không cần retraining model.
    \item \textbf{Transparent} – Có thể theo dõi nguồn thông tin.
    \item \textbf{Specificity} – Câu trả lời phù hợp domain (y tế, pháp luật, v.v.).
\end{itemize}

\subsubsection{Graph Retrieval Augmented Generation (GraphRAG)}

\textbf{Nâng cấp RAG:} Kết hợp \textbf{Knowledge Graph} (đồ thị tri thức) với Vector Store.

\begin{itemize}
    \item \textbf{Knowledge Graph (Neo4j)} – Biểu diễn mối quan hệ có cấu trúc:
    \begin{verbatim}
(Symptom: "Đau dạ dày")-[:INDICATES]->(Disease: "Gastric Ulcer")
(Disease: "Gastric Ulcer")-[:TREATED_BY]->(Product: "Omeprazole")
(Product: "Omeprazole")-[:CATEGORY]->(Category: "PPI")
(User: "Patient123")-[:SEARCHED]->(Symptom: "Đau dạ dày")
    \end{verbatim}
    Lợi ích: Tìm kiếm theo logic quan hệ, không chỉ semantic similarity.
    
    \item \textbf{Vector Store (FAISS)} – Tìm kiếm phi cấu trúc dựa mức độ tương đồng vector:
    \begin{verbatim}
Query: "Tôi bị đau bụng vào buổi sáng sớm"
→ Vector encode → FAISS search top-5 similar chunks
→ Chunks về GERD, Peptic Ulcer, IBS, ...
    \end{verbatim}
    Lợi ích: Semantic search, catch queries không match exact keywords.
\end{itemize}

\textbf{Combine 3 nguồn:}
\begin{enumerate}
    \item Neo4j Knowledge Graph – Truy xuất facts có cấu trúc.
    \item FAISS Vector Store – Truy xuất docs phi cấu trúc.
    \item Product Service REST API – Fetch thông tin sản phẩm thực tế.
\end{enumerate}

\subsubsection{Phân Tích Hành Vi Khách Hàng (Customer Behavior Analysis)}

\textbf{Mục đích:} Hiểu nhu cầu, sở thích của khách hàng để gợi ý phù hợp.

\textbf{Kỹ thuật:}
\begin{itemize}
    \item \textbf{Collaborative Filtering} – Assume: người dùng có hành vi tương tự → nên like sản phẩm tương tự.
    \[ \text{similarity}(u, v) = \frac{\vec{u} \cdot \vec{v}}{|\vec{u}| |\vec{v}|} \]
    
    \item \textbf{Content-Based Filtering} – Recommend dựa trên sản phẩm mà user đã like.
    
    \item \textbf{Hybrid} – Kết hợp cả hai.
\end{itemize}

\textbf{Trong Health-Micro-AI:}
\begin{itemize}
    \item Ghi nhận hành vi: product viewed, products searched, products purchased.
    \item Neo4j node: (User)-[:VIEWED {weight}]->(Product), (User)-[:SEARCHED]->(Symptom).
    \item Query: "Những người xem sản phẩm X cũng xem sản phẩm Y" → recommend Y.
\end{itemize}

\subsubsection{Graph Neural Networks (GNN) và SPD Manifold}

\textbf{Graph Neural Networks:} Mô hình học sâu tác động trên đồ thị. Mỗi node học biểu diễn (embedding) kết hợp thông tin từ neighbors.

\begin{itemize}
    \item \textbf{GraphSAGE} – Sample neighbors, aggregate their embeddings.
    \item \textbf{GAT} – Attention mechanism, trọng số cạnh động.
    \item \textbf{GCN} – Graph Convolutional Networks.
\end{itemize}

\textbf{SPD Manifold} (Symmetric Positive Definite): Không gian toán học đặc biệt dùng để bảo toàn cấu trúc không gian khi learn embeddings. Giúp recommendation không bị "collapse" vào vài chiều.

\textbf{Ứng dụng:} Dự đoán liên kết (link prediction) – "Người dùng này nên được recommend sản phẩm nào?" dựa trên cấu trúc đồ thị.

\subsubsection{Học Sâu (Deep Learning)}

\textbf{Định nghĩa:} Mạng neural nhiều lớp học biểu diễn phân cấp từ dữ liệu thô.

\textbf{Kiến trúc phổ biến:}
\begin{itemize}
    \item \textbf{Convolutional Neural Networks (CNN)} – Xử lý ảnh, temporal data.
    \item \textbf{Recurrent Neural Networks (RNN/LSTM)} – Xử lý sequence, time-series.
    \item \textbf{Transformer} – Self-attention, machine translation, LLM.
    \item \textbf{Graph Neural Networks (GNN)} – Xử lý đồ thị.
\end{itemize}

\textbf{Trong project:} Không train Deep Learning model locally. Thay vào đó:
\begin{itemize}
    \item Dùng Gemini API (pre-trained LLM) để embed text và sinh text.
    \item Dùng pre-trained model (SentenceTransformer) để embed medical documents.
    \item GNN dùng thư viện DGL hoặc PyTorch Geometric (optional, cho production).
\end{itemize}

\subsubsection{Cơ Sở Tri Thức (Knowledge Base)}

\textbf{Định nghĩa:} Tập hợp thông tin khoa học được cấu trúc để các hệ thống AI có thể query.

\textbf{Trong Health-Micro-AI:}
\begin{itemize}
    \item \textbf{Neo4j Knowledge Base} – Các mối quan hệ:
    \begin{verbatim}
Symptom --- INDICATES --> Disease --- TREATED_BY --> Product
User --- SEARCHED --> Symptom
User --- PURCHASED --> Product
Disease --- DIAGNOSED_BY --> Diagnostic
    \end{verbatim}
    
    \item \textbf{Vector Knowledge Base (FAISS)} – Các tài liệu y khoa được tokenize và embed:
    \begin{verbatim}
Chunks từ: benhdaday.csv, faq.txt, medical literature
→ Embed bằng Gemini text-embedding API
→ Lưu vào FAISS Index
→ Tìm kiếm bằng cosine similarity
    \end{verbatim}
\end{itemize}

\subsection{Ứng Dụng Thực Tế vào E-commerce Healthcare}

\subsubsection{Clinical-Advisory Service: Chatbot GraphRAG}

\textbf{Chức năng:} Trả lời câu hỏi y tế, gợi ý sản phẩm, cảnh báo.

\textbf{Luồng xử lý:}

\begin{enumerate}
    \item \textbf{Intent Classification} – Phân loại câu hỏi: triệu chứng, chẩn đoán, thuốc, etc.
    
    \item \textbf{Entity Extraction} – Rút trích entities (bệnh, triệu chứng, sản phẩm) từ query.
    
    \item \textbf{Neo4j Graph Traversal} – Query knowledge graph:
    \begin{verbatim}
MATCH (s:Symptom{name: "Đau dạ dày"})
      -[r:INDICATES]->
      (d:Disease)
      -[s:TREATED_BY]->
      (p:Product)
RETURN d.name, p.name, p.product_id
    \end{verbatim}
    Trong ví dụ trên, tìm mối quan hệ từ "triệu chứng đau dạ dày" → bệnh → sản phẩm.
    
    \item \textbf{FAISS Vector Search} – Tìm top-5 docs y khoa tương đồng:
    \begin{verbatim}
query_vector = embed("Tôi bị đau dạ dày phải làm sao?")
top_5_chunks = faiss_index.search(query_vector, k=5)
# Kết quả: [chunk_về_peptic_ulcer, chunk_về_GERD, ...]
    \end{verbatim}
    
    \item \textbf{Context Assembly} – Ghép Graph results + Vector results + Product API data:
    \begin{verbatim}
context = """
Kiến thức từ Graph:
- Triệu chứng 'Đau dạ dày' có thể chỉ định các bệnh:
  - Peptic Ulcer (điều trị bằng Omeprazole, Esomeprazole)
  - GERD (điều trị bằng PPI)

Kiến thức từ Docs:
- [Peptic Ulcer Disease]: Là tổn thương niêm mạc dạ dày...
- [GERD]: Trào ngược acid có thể gây...

Sản phẩm hiện có:
- Omeprazole 20mg (giá 50k, trong kho)
- Esomeprazole 40mg (giá 150k, trong kho)
"""
    \end{verbatim}
    
    \item \textbf{Gemini LLM Generation} – Yêu cầu Gemini sinh câu trả lời từ context:
    \begin{verbatim}
prompt = f"""
Bạn là bác sĩ tư vấn y tế chuyên nghiệp.
Dựa trên thông tin nền tảng:
{context}

Câu hỏi của khách hàng: "Tôi bị đau dạ dày, có nên dùng Omeprazole không?"

Trả lời một cách khoa học, an toàn, có trích dẫn.
Nếu cần, gợi ý ghé bác sĩ.
"""
response = genai.generate_text(prompt)
    \end{verbatim}
    
    \item \textbf{Medical Disclaimer} – Thêm tuyên bố từ chối trách nhiệm:
    \begin{verbatim}
DISCLAIMER: "Thông tin này chỉ mang tính chất tham khảo. 
Vui lòng tham khảo ý kiến bác sĩ chuyên khoa trước khi 
sử dụng bất kỳ sản phẩm nào."
    \end{verbatim}
\end{enumerate}

\textbf{Ưu điểm và Sự Tích Hợp E-Commerce Xuyên Suốt:}
\begin{itemize}
    \item \textbf{Grounded và Safe} – Câu trả lời dựa trên knowledge y khoa đã được xác thực, kèm theo disclaimer để đảm bảo an toàn.
    \item \textbf{Tích Hợp E-Commerce Chặt Chẽ} – Không giống như các chatbot thông thường chỉ trả lời chữ, AI của hệ thống nhận diện insight và \textbf{gắn trực tiếp Data Sản Phẩm thực tế} (như Omeprazole 20mg, giá 35.000 VND).
    \item \textbf{Từ Tư Vấn Đến Chốt Đơn (Conversion)} – Các gợi ý thuốc được AI sinh ra dưới dạng Link hoặc Card sản phẩm. Bệnh nhân có thể ngay lập tức click vào "Thêm vào giỏ hàng" (Add to Cart) để chuyển đổi từ người đi hỏi thành khách mua hàng. Đây là lõi sức mạnh của \textit{E-Commerce AI Integration}.
\end{itemize}

\subsubsection{Recommender-AI Service: Collaborative Filtering}

\textbf{Chức năng:} Gợi ý sản phẩm dựa trên hành vi người dùng tương tự.

\textbf{Thuật toán Jaccard Similarity:}
$$J(u, v) = \frac{|I_u \cap I_v|}{|I_u \cup I_v|}$$

Trong đó $I_u$, $I_v$ là tập sản phẩm mà user $u$, $v$ đã view/purchase.

\textbf{Luồng:}
\begin{enumerate}
    \item Xây dựng user-item matrix: mỗi hàng là user, cột là product, giá trị là rating/weight.
    
    \item Tính similarity giữa current user và tất cả users khác dùng Jaccard.
    
    \item Lấy top-N similar users.
    
    \item Gợi ý products mà similar users viewed/purchased nhưng current user chưa.
    
    \item Rank theo confidence score (số similar users → product).
\end{enumerate}

\textbf{Pseudocode:}
\begin{verbatim}
def collaborative_filtering_recommend(user_id, k=5, n=10):
    user_items = get_user_items(user_id)  # set of product IDs
    
    similar_users = []
    for other_user in all_users:
        if other_user == user_id:
            continue
        other_items = get_user_items(other_user)
        similarity = jaccard_similarity(user_items, other_items)
        if similarity > 0:
            similar_users.append((other_user, similarity))
    
    similar_users.sort(by=similarity, desc=True)
    top_k_similar = similar_users[:k]
    
    candidates = {}  # product_id -> score
    for other_user, sim_score in top_k_similar:
        for product in get_user_items(other_user):
            if product not in user_items:
                candidates[product] += sim_score
    
    top_n_recommend = sorted(candidates.items(), 
                             by=score, desc=True)[:n]
    return top_n_recommend
\end{verbatim}

\textbf{Ưu điểm:}
\begin{itemize}
    \item \textbf{Đơn giản} – Dễ implement, không cần training.
    \item \textbf{Nhanh} – Compute on-the-fly.
    \item \textbf{Explainable} – Có thể nói: "Những người khác giống bạn cũng thích sản phẩm X".
\end{itemize}

\textbf{Nhược điểm:}
\begin{itemize}
    \item \textbf{Cold start} – User mới chưa có hành vi.
    \item \textbf{Sparse data} – Nếu users ít, khó tính similarity.
\end{itemize}

% ══════════════════════════════════════════════════════════════════════════════
% CHƯƠNG 5: TRIỂN KHAI VÀ MÃ NGUỒN
% ══════════════════════════════════════════════════════════════════════════════

\newpage
\section{Triển Khai và Mã Nguồn}

\subsection{Cấu Trúc Thư Mục}

\begin{verbatim}
health-micro-ai/
+-- docker-compose.yml
+-- build.bat / build.sh
+-- setup_db.bat
+-- seed_all.py
+-- init_postgres.sql
+-- .env.example
+-- api-gateway/
    +-- manage.py, requirements.txt
    +-- app/
    |   +-- views.py, middleware.py, models.py
    +-- templates/ (index.html, ai_chat.html, login.html)
    +-- Dockerfile
+-- auth-service/
    +-- main.py, domain/, application/, infrastructure/
    +-- Dockerfile
+-- patient-service/
    +-- domain/, application/, infrastructure/
    +-- Dockerfile
+-- pharmacy-service/
    +-- domain/, interfaces/, infrastructure/
    +-- Dockerfile
+-- prescription-service/
    +-- domain/, application/
    +-- Dockerfile
+-- clinical-advisory-service/
    +-- main.py, graph_rag/, domain/, services/
    +-- Dockerfile
+-- recommender-ai-service/
    +-- app.py, domain/
    +-- Dockerfile
+-- shared_ddd/
    +-- base.py, event_bus.py, event_bus_integration.py
+-- README.md, API_QUICK_TEST.md, BAO_CAO_TIEU_LUAN.tex
\end{verbatim}

\subsection{Các Hàm Quan Trọng (Core Functions)}

\subsubsection{1. GraphRAG Query với Neo4j + FAISS}

\begin{lstlisting}[language=Python, caption=RAG Engine Core (clinical-advisory-service/graph\_rag/rag\_engine.py)]
import faiss
import numpy as np
from neo4j import GraphDatabase
import google.generativeai as genai

class GraphRAGEngine:
    def __init__(self, neo4j_uri, neo4j_user, neo4j_password, 
                 gemini_api_key, faiss_index_path):
        self.driver = GraphDatabase.driver(neo4j_uri, 
                                           auth=(neo4j_user, neo4j_password))
        genai.configure(api_key=gemini_api_key)
        self.model = genai.GenerativeModel('gemini-1.5-flash')
        
        # Load FAISS index
        self.faiss_index = faiss.read_index(
            f"{faiss_index_path}/medical.index")
        with open(f"{faiss_index_path}/chunks.json") as f:
            self.chunks = json.load(f)
    
    def query_graph(self, symptom: str) -> List[Dict]:
        """Query Neo4j for symptom -> disease -> product relationships"""
        with self.driver.session() as session:
            result = session.run("""
                MATCH (s:Symptom{name: $symptom})
                      -[r:INDICATES]->
                      (d:Disease)
                      -[t:TREATED_BY]->
                      (p:Product)
                RETURN d.name as disease, 
                       d.icd_code, 
                       p.product_id,
                       p.name as product_name,
                       t.priority
                ORDER BY t.priority DESC
                LIMIT 10
            """, symptom=symptom)
            return [dict(record) for record in result]
    
    def search_vectors(self, query: str, k: int = 5) -> List[str]:
        """Search similar medical documents using FAISS"""
        # Embed query using Gemini
        response = genai.embed_content(
            model="models/embedding-001",
            content=query,
            task_type="RETRIEVAL_QUERY"
        )
        query_vector = np.array(response['embedding'], dtype=np.float32)
        query_vector = query_vector / np.linalg.norm(query_vector)
        
        # FAISS search
        distances, indices = self.faiss_index.search(
            np.array([query_vector]), k=k)
        
        results = []
        for idx in indices[0]:
            if idx < len(self.chunks):
                results.append(self.chunks[idx])
        return results
    
    def generate_answer(self, question: str, 
                       graph_context: List[Dict],
                       vector_context: List[str]) -> str:
        """Using Gemini to synthesize answer from context"""
        # Build context
        context = f"""
You are a healthcare advisor. Based on the following information:

GRAPH KNOWLEDGE:
{json.dumps(graph_context, ensure_ascii=False, indent=2)}

MEDICAL DOCUMENTS:
{chr(10).join(vector_context)}

USER QUESTION: {question}

Provide a helpful, evidence-based answer. 
Include products that may help.
Always add a medical disclaimer.
"""
        
        response = self.model.generate_content(context)
        return response.text
\end{lstlisting}

\subsubsection{2. Collaborative Filtering Recommendation}

\begin{lstlisting}[language=Python, caption=Collaborative Filtering (recommender-ai-service/app.py)]
from collections import defaultdict
import json

class CollaborativeFilteringEngine:
    def __init__(self):
        self.user_items = {}  # user_id -> set of product_ids
    
    def jaccard_similarity(self, set1, set2):
        """Calculate Jaccard similarity between two sets"""
        if not set1 or not set2:
            return 0.0
        intersection = len(set1 & set2)
        union = len(set1 | set2)
        return intersection / union if union > 0 else 0.0
    
    def get_recommendations(self, user_id, 
                           k_similar_users=5, 
                           n_recommendations=10):
        """
        Recommend products to user based on similar users
        
        Args:
            user_id: Current user
            k_similar_users: Top-K similar users
            n_recommendations: Number of products to recommend
            
        Returns:
            List of (product_id, confidence_score) tuples
        """
        
        if user_id not in self.user_items:
            return []  # No history for this user
        
        user_items = self.user_items[user_id]
        
        # Calculate similarity with all other users
        similarities = []
        for other_user, other_items in self.user_items.items():
            if other_user == user_id:
                continue
            
            sim = self.jaccard_similarity(user_items, 
                                         set(other_items))
            if sim > 0:
                similarities.append((other_user, sim))
        
        # Get top-K similar users
        similarities.sort(key=lambda x: x[1], reverse=True)
        top_similar = similarities[:k_similar_users]
        
        # Aggregate candidate products
        candidate_scores = defaultdict(float)
        for similar_user, sim_score in top_similar:
            for product in self.user_items.get(similar_user, []):
                if product not in user_items:
                    candidate_scores[product] += sim_score
        
        # Sort and return top-N
        recommendations = sorted(candidate_scores.items(),
                                key=lambda x: x[1], 
                                reverse=True)[:n_recommendations]
        
        return recommendations
    
    def record_interaction(self, user_id, product_id, 
                          interaction_type="view"):
        """Record user-product interaction"""
        if user_id not in self.user_items:
            self.user_items[user_id] = []
        
        if product_id not in self.user_items[user_id]:
            self.user_items[user_id].append(product_id)
\end{lstlisting}

\subsubsection{3. API Gateway Routing}

\begin{lstlisting}[language=Python, caption=API Gateway Routing (api-gateway/app/views.py)]
from django.http import JsonResponse
from django.views import View
import httpx
import json

class GatewayProxyView(View):
    """Proxy requests to microservices"""
    
    SERVICE_REGISTRY = {
        'auth': 'http://auth-service:8000',
        'patient': 'http://patient-service:8000',
        'pharmacy': 'http://pharmacy-service:8000',
        'prescription': 'http://prescription-service:8000',
        'clinical-advisory': 'http://clinical-advisory-service:8000',
        'recommender': 'http://recommender-ai-service:8000',
    }
    
    def route_request(self, request, service, path):
        """Route HTTP request to appropriate service"""
        
        # Verify JWT token
        auth_header = request.headers.get('Authorization', '')
        if not auth_header.startswith('Bearer '):
            return JsonResponse(
                {'error': 'Unauthorized'}, 
                status=401
            )
        
        # Get service URL
        if service not in self.SERVICE_REGISTRY:
            return JsonResponse(
                {'error': f'Unknown service: {service}'}, 
                status=404
            )
        
        service_url = self.SERVICE_REGISTRY[service]
        
        # Build target URL
        target_url = f"{service_url}/api/{path}"
        
        # Forward request using httpx
        try:
            with httpx.Client(timeout=10.0) as client:
                # Copy headers, preserve JWT
                headers = dict(request.headers)
                
                if request.method == 'GET':
                    response = client.get(target_url, headers=headers)
                elif request.method == 'POST':
                    response = client.post(
                        target_url, 
                        headers=headers,
                        content=request.body
                    )
                elif request.method == 'PUT':
                    response = client.put(
                        target_url, 
                        headers=headers,
                        content=request.body
                    )
                else:
                    return JsonResponse(
                        {'error': 'Method not allowed'}, 
                        status=405
                    )
                
                # Return upstream response
                return JsonResponse(
                    response.json(), 
                    status=response.status_code
                )
        
        except httpx.TimeoutException:
            return JsonResponse(
                {'error': 'Service timeout'}, 
                status=504
            )
        except Exception as e:
            return JsonResponse(
                {'error': str(e)}, 
                status=500
            )
\end{lstlisting}

\subsection{Giao Diện Hệ Thống (Screenshots)}

Dưới đây là các placeholder cho hình ảnh giao diện web:

\begin{figure}[H]
    \centering
    \includegraphics[width=0.8\textwidth]{screenshots/homepage.png}
    \caption{Trang chủ – Giao diện người dùng bệnh nhân}
    \label{fig:homepage}
\end{figure}

\begin{figure}[H]
    \centering
    \includegraphics[width=0.8\textwidth]{screenshots/ai_chatbot.png}
    \caption{UI Chatbot AI – Tư vấn y tế theo GraphRAG}
    \label{fig:chatbot}
\end{figure}

\begin{figure}[H]
    \centering
    \includegraphics[width=0.8\textwidth]{screenshots/product_search.png}
    \caption{Tìm kiếm sản phẩm – Với filter theo triệu chứng}
    \label{fig:search}
\end{figure}

\begin{figure}[H]
    \centering
    \includegraphics[width=0.8\textwidth]{screenshots/prescription_management.png}
    \caption{Quản lý đơn thuốc – Tạo, edit, submit}
    \label{fig:prescription}
\end{figure}

\begin{figure}[H]
    \centering
    \includegraphics[width=0.8\textwidth]{screenshots/recommendations.png}
    \caption{Gợi ý cá nhân hóa – Dựa trên collaborative filtering}
    \label{fig:recommendations}
\end{figure}

% ══════════════════════════════════════════════════════════════════════════════
% KẾT LUẬN
% ══════════════════════════════════════════════════════════════════════════════

\newpage
\section{Kết Luận}

Dự án \textbf{Health-Micro-AI} là một thực hiện toàn diện của kiến trúc Microservices kết hợp với Domain-Driven Design và Trí tuệ Nhân tạo. Thông qua việc tách ứng dụng thành 11 services độc lập, mỗi service quản lý một Bounded Context, hệ thống đạt được:

\begin{enumerate}
    \item \textbf{Khả năng mở rộng cao} – Từng service có thể scale riêng theo load.
    
    \item \textbf{Tính linh hoạt} – Mỗi team phát triển một service mà không ảnh hưởng service khác.
    
    \item \textbf{Khả năng phục hồi} – Lỗi một service không làm crash toàn hệ.
    
    \item \textbf{Trí tuệ Nhân tạo tích hợp} – Clinical-Advisory Service cung cấp tư vấn y tế dựa GraphRAG, giúp người dùng ra quyết định tốt hơn.
    
    \item \textbf{Gợi ý cá nhân hóa} – Recommender Services cung cấp sản phẩm phù hợp cho từng người dùng.
\end{enumerate}

\textbf{Những thách thức được giải quyết:}
\begin{itemize}
    \item Distributed transaction – Sử dụng event-driven architecture qua RabbitMQ.
    \item Data consistency – Eventual consistency model, acceptable cho healthcare.
    \item Monitoring & Logging – Correlation IDs, centralized logging.
\end{itemize}

\textbf{Hướng phát triển tương lai:}
\begin{itemize}
    \item Thêm API chính thức cho third-party integrations.
    \item Implement Kubernetes orchestration cho production deployment.
    \item Advanced ML models: GNN training on Neo4j graphs, SPD manifold.
    \item Multi-language support: Mở rộng AI tư vấn cho tiếng Anh, ngôn ngữ khác.
    \item Mobile app: Native iOS/Android client.
    \item Analytics dashboard: Real-time insights về dữ liệu sức khỏe thatuser.
\end{itemize}

\textbf{Kết luận:} Kiến trúc này chứng minh rằng Microservices + DDD + AI là combination mạnh mẽ cho các ứng dụng healthcare phức tạp. Với sự phân tách rõ ràng giữa các domains và công nghệ AI hiện đại, hệ thống vừa dễ bảo trì, vừa mang lại giá trị cao cho người dùng.

\end{document}
